\section{\texorpdfstring{$E_8$}{E8} from nine histories}
\label{sec:e8}

$E_8$ is the largest of the five exceptional simple Lie algebras: $248$ dimensions, $240$
roots, a symmetry so rigid that it appears, often unexpectedly, in string theory, in the
densest packing of spheres in eight dimensions, and in the physics of certain magnetic
materials. This section shows that one qutrit's nine operator histories build $E_8$ in a
way in which every piece has a meaning in the finite geometry --- and it records
carefully one famous coincidence that turned out to be a trap.

\begin{plainwords}
A Lie algebra is the ``infinitesimal'' version of a continuous symmetry: the set of all
directions in which you can nudge a system while keeping its rules intact, together with
the rule (the \emph{bracket}) for what happens when two nudges are combined in opposite
orders. Its \emph{roots} are the special directions that the symmetry's own ``clock''
(its maximal commuting part) rotates at definite rates. $E_8$ has $240$ of them.
\end{plainwords}

\subsection{Eighty plus eighty-four plus eighty-four}

Recall from Section~3 that the $80$ non-identity two-sided Weyl maps on one qutrit span
$\sll_9$, the Lie algebra of traceless $9\times9$ matrices acting on the nine operator
histories $\ket{i}\!\bra{j}$. Now add \emph{packets}: unordered triples of histories. There
are $\binom93=84$ of them, forming the representation $\Lambda^3\C^9$; their duals form
$\Lambda^6\C^9$, another $84$. A classical construction \cite{Adams1996,Baez2002} makes
\[
E_8 \;=\; \sll_9\;\oplus\;\Lambda^3\C^9\;\oplus\;\Lambda^6\C^9,\qquad
248=80+84+84 ,
\]
into a Lie algebra. On the nine history cells the brackets are literal packet operations:
$E_{ij}$ replaces cell $j$ by cell $i$ inside a packet; two disjoint triples bracket to the
complementary six-set; a triple and a six-set differing by one cell contract to an
$E_{ij}$; opposite packets return a Cartan direction. \tierC\
(\cert{w33\_20260924\_temporal\_su9\_e8\_compiler.py},
\cert{w33\_20260924\_temporal\_a8\_e8\_process\_bracket.py})

\subsection{Hesse lines are copies of \texorpdfstring{$A_2$}{A2}}

Place the nine histories on the affine plane $\F_3^2$ as in Figure~\ref{fig:hesse}. A
triple $S$ of cells labels the root
\[
w_S=\sum_{i\in S}e_i-\tfrac13\sum_{i=1}^9e_i ,\qquad w_S\cdot w_T=|S\cap T|-1 .
\]
(Indeed $|w_S|^2=3-2+1=2$.) The $84$ triples fall into two orbits of the affine group
$AGL(2,3)$ (order $432$): the $12$ \emph{lines} of the plane and the $72$ non-collinear
\emph{triangles}.

\begin{result}{Theorem 5.1 (four clocks, four $A_2$'s) \hfill \tierT\ \tierC}
The three lines of one parallel class give three roots with pairwise products $-1$ and
sum $0$; with their negatives they form a root system $A_2$. Lines of different classes
meet in exactly one cell, so their roots are orthogonal. The $12$ Hesse lines therefore
span a reflection-closed subsystem $4A_2\subset E_8$, one $A_2$ per clock, and the root
shell refines as
\[
240\;=\;\underbrace{72}_{A_8}\;+\;\underbrace{24}_{\text{Hesse }4A_2}\;+\;\underbrace{144}_{\text{triangles}} .
\]
\cert{w33\_20260924\_temporal\_hesse\_4a2\_a8\_bridge.py}
\end{result}

\begin{proof}
Disjoint triples have $|S\cap T|=0$, hence product $-1$; the three lines of a class
partition the nine cells, so $\sum w_S=\sum_i e_i-3\cdot\frac13\sum_ie_i=0$. Two
non-parallel lines of an affine plane meet in one point, giving product $1-1=0$.
\end{proof}

\begin{figure}[htbp]
\centering
\begin{tikzpicture}[scale=0.95]
  % four A2 hexagons, one per clock
  \foreach \c/\n/\x in {teal/1/0,sky/2/2.6,gold/3/5.2,sage/4/7.8}{
    \begin{scope}[xshift=\x cm]
      \foreach \k in {0,...,5}{
        \draw[\c,line width=0.9pt,-{Stealth[length=1.6mm]}] (0,0) -- (60*\k+30:0.9);}
      \draw[\c!35,line width=0.5pt] (30:0.9) \foreach \k in {1,...,5}{ -- (60*\k+30:0.9)} -- cycle;
      \fill[\c] (0,0) circle (1.3pt);
      \node[font=\sffamily\scriptsize\bfseries,text=\c!80!black] at (0,-1.25) {$A_2$, clock \n};
    \end{scope}}
  \node[font=\sffamily\scriptsize,text=slate] at (3.9,-1.85) {$4\times 6=24$ roots: the Hesse subsystem};
  % glue sectors
  \begin{scope}[xshift=11.3cm,yshift=0cm]
    \foreach \k in {0,...,7}{
      \pgfmathsetmacro{\a}{\k*45+22.5}
      \fill[rose!18,draw=rose!60] (\a:1.05) circle (0.33);
      \node[font=\tiny] at (\a:1.05) {$27$};}
    \node[font=\sffamily\scriptsize,align=center] at (0,0) {tetracode\\glue};
    \node[font=\sffamily\scriptsize,text=slate,align=center] at (0,-1.85) {$8$ nonzero codewords\\$\times\,27=216$ roots};
  \end{scope}
\end{tikzpicture}
\caption{$E_8$ assembled from four clocks. Each Hesse clock contributes one $A_2$ root
system (left). The remaining $216$ roots come in eight sectors of $27$, one for each
nonzero word of the tetracode (right): $240=4\cdot6+8\cdot27$.}
\label{fig:a24}
\end{figure}

\subsection{The glue is the tetracode, and the tetracode is the clock}

The lattice $E_8$ contains the sublattice $A_2^4$ with quotient $\F_3^2$, and the
well-known construction of $E_8$ from the tetracode \cite{ConwaySloane} reads it as
\[
E_8=\bigcup_{c\in\mathcal{T}}\bigl(c+A_2^{\,4}\bigr),\qquad
240=4\cdot6+8\cdot27 ,
\]
where $\mathcal{T}\subset\F_3^4$ is the tetracode: each of the eight nonzero codewords has
weight three, and contributes $3^3=27$ minimal vectors. Combined with Section~4, the
four glue coordinates are the four clocks and the eight nonzero sectors are the eight
oriented, nonzero linear clock readings. \tierT\ (classical) \ \tierC\ (the clock reading,
\cert{w33\_pass10946\_clock\_code\_cone\_objectwise.py})

\subsection{A trap: \texorpdfstring{$240=240$}{240 = 240}}
\label{sec:e8-trap}

The collinearity graph of $\W$ has $240$ edges, and $E_8$ has $240$ roots. For years the
repository looked for a natural bijection between them. The question has a precise and
negative answer.

\begin{result}{Theorem 5.2 (no equivariant bijection) \hfill \tierC}
The two groups of order $51\,840$ of Section~2 each do only half the job:
\begin{itemize}
\item $\PGSp(4,3)\cong W(E_6)$ acts faithfully and transitively on the $240$ edges of $\W$,
      but \emph{intransitively} on the $240$ roots of $E_8$;
\item $\Sp(4,3)$ acts faithfully and transitively on the $240$ roots, but its centre acts
      trivially on the edges, so its edge action is not faithful.
\end{itemize}
Hence no bijection between edges and roots is equivariant for a common group.
\cert{w33\_pass1020\_e8\_transitive\_51840.g}
\end{result}

What \emph{does} exist is better. The Weyl group of $E_8$ contains a regular element $w$
of order three. By Springer's theory \cite{Springer} its centraliser is the complex reflection group
$G_{32}\cong\Z_3\times\Sp(4,3)$ of order $155\,520$ --- the symmetry group of the
\emph{Witting polytope}, whose $240$ vertices in $\C^4$ are the $E_8$ roots.

\begin{result}{Theorem 5.3 (the $6:1$ fibration) \hfill \tierC}
The $240$ roots of $E_8$ fall into $40$ complex lines $\{\pm v,\pm\omega v,\pm\omega^2 v\}$,
each a copy of the six Eisenstein units. $\Sp(4,3)$ permutes the $40$ lines with
subdegrees $1,12,27$: the quotient is the collinearity graph of $\W$ --- on its
\emph{points}, not its lines. The fibre group is $\langle c^5\rangle\cong\Z_6$ for a
Coxeter element $c$ of $W(E_8)$.
\cert{w33\_pass1021\_e8\_fibration\_over\_forty.g}
\end{result}

\begin{figure}[htbp]
\centering
\begin{tikzpicture}[scale=1]
  % base: 40 points (sketch: the 1+12+27 view)
  \begin{scope}
    \fill[ink] (0,0) circle (2.2pt);
    \foreach \k in {0,...,11}{\fill[teal] ({30*\k}:0.9) circle (1.6pt);}
    \foreach \k in {0,...,26}{\fill[rose!70] ({360*\k/27+5}:1.6) circle (1.3pt);}
    \node[font=\sffamily\scriptsize,text=slate] at (0,-2.15) {base: the $40$ points of $\W$};
  \end{scope}
  % fibre over one point
  \draw[-{Stealth[length=2.5mm]},ink!60,line width=0.9pt] (3.1,1.3) to[bend right=25] node[above,font=\sffamily\scriptsize,pos=0.35]{project} (0.12,0.12);
  \begin{scope}[xshift=4.2cm,yshift=1.5cm]
    \foreach \k/\lab in {0/{v},1/{-\omega^2v},2/{\omega v},3/{-v},4/{\omega^2 v},5/{-\omega v}}{
      \fill[gold] (60*\k:1.0) circle (2.3pt);
      \node[font=\scriptsize] at (60*\k:1.42) {$\lab$};}
    \draw[gold!60] (0:1) \foreach \k in {1,...,5}{-- (60*\k:1)} -- cycle;
    \node[font=\sffamily\scriptsize,text=slate,align=center] at (0,-1.85) {fibre: $6$ roots $=$ one\\Witting diameter, $\Z_6$};
  \end{scope}
  \node[font=\sffamily\small,align=left,anchor=west] at (6.4,0.4)
   {$240\;=\;40\times6$\\[3pt]\scriptsize roots $\to$ points,\\\scriptsize $\Sp(4,3)$-equivariant};
\end{tikzpicture}
\caption{The correct relation between $E_8$ and $\W$: not edges to roots, but a
six-to-one map from roots to points, whose fibres are the Eisenstein unit hexagons.}
\label{fig:fibration}
\end{figure}

\begin{figure}[p]
\centering
\begin{tikzpicture}[scale=1.12]
  \input{fig_e8coxeter}
  \fill[ink] (0,0) circle (1.2pt);
\end{tikzpicture}
\caption{The $240$ roots of $E_8$ projected to the Coxeter plane of a Coxeter element $c$
(order $30$), which acts on the plane as a rotation by $12^\circ$. The roots lie on eight
rings of thirty. The orbits of $c^5$ --- rotation by $60^\circ$ --- are regular hexagons,
five per ring: these $40$ hexagons are the $40$ fibres of Theorem~5.3, one for each point of
$\W$. Each hexagon is $\{\pm v,\pm\omega v,\pm\omega^2v\}$, since $c^{15}=-1$ and $c^{10}$ has
order three. All coordinates computed (\texttt{make\_figures.py}).}
\label{fig:coxeter}
\end{figure}

\begin{physical}[What it means physically: a photon's four-dimensional state space]
The $40$ Witting diameters are $40$ rays in $\C^4$ --- the state space of a photon's path
and polarisation, or of two qubits. Two rays are orthogonal exactly when the corresponding
points of $\W$ are collinear. So the same geometry is drawn twice: by \emph{states} of a
four-level system (the Witting rays) and by \emph{operators} on two qutrits (the Pauli
classes). The photonic Holonet papers call this operator--operand duality.
\tierC\ (\cert{bt817\_self\_entangled\_photon\_atlas.py}, \cert{bt821\_operator\_operand\_duality.py})
\end{physical}

\subsection{\texorpdfstring{$E_6\times A_2$}{E6 x A2}, the twenty-seven, and a cubic}

Under its subalgebra $E_6\oplus A_2$,
\[
E_8=(\mathbf{78},\mathbf1)\oplus(\mathbf1,\mathbf8)\oplus(\mathbf{27},\mathbf3)\oplus(\overline{\mathbf{27}},\overline{\mathbf3}),
\]
which is a $\Z_3$-grading $248=86+81+81$. The $27$ of $E_6$ carries a unique invariant
cubic form $N$; in the repository's committed basis it has exactly $45$ signed terms
(\emph{triads}), each coordinate in $5$ triads and each pair in at most one ---
precisely the incidence of the $27$ lines and $45$ tritangent planes of a smooth cubic
surface, whose symmetry group is $W(E_6)$, the automorphism group of $\W$ \cite{Garibaldi}.
\tierT\ (classical) \ \tierC\ (\cert{canonical\_su3\_gauge\_and\_cubic.json})

\begin{physical}
In grand unified theories built on $E_6$, the $\mathbf{27}$ holds one generation of
matter: under $SO(10)$ it splits $\mathbf{27}=\mathbf{16}+\mathbf{10}+\mathbf1$ ---
sixteen fermions (including a right-handed neutrino), a ten of Higgs-type fields, and a
singlet. The cubic $N$ is then the unique renormalisable Yukawa coupling
$\mathbf{27}^3$. Section~8 shows this same split arising as scalar, vector and spinor of a
ten-dimensional Lorentz group.
\end{physical}

\subsection{The explicit semisimple three-plane and its qutrit mirrors}

The $\Z_3$ grading has rank three in Vinberg's sense, but that abstract statement does
not identify a Cartan subspace in the repository's signed $3\otimes27$ coordinates. An
earlier displayed commuting kernel was nilpotent and its Cartan interpretation was
withdrawn. The missing object can now be written and checked exactly.

\begin{result}{Theorem 5.4 (the signed cubic contains its $G_{26}$ Cartan plane)
\hfill \tierC}
There is an integral signed-cubic background $v$ whose cubic Jacobian has rank $78$ and
whose grade-one centraliser is a three-dimensional commuting plane $\mathfrak c$. For
$N=18\,\mathrm{ad}(v)$,
\[
 \operatorname{rank} N=\operatorname{rank} N^3=234.
\]
The three diagonal blocks of $N^3$ have rank $78$ and the same nonzero characteristic
polynomial. The only repeated factors have degree nine and each has its full
$27$-dimensional kernel, so $N^3$, hence $N$, is diagonalizable. Thus $\mathfrak c$ is a
regular semisimple Cartan subspace. \cert{w33\_pass11260\_exact\_semisimple\_cartan.py}

If $(x,y,z)$ are its committed rational coordinates, an explicit invertible map over
$\Q(\sqrt[3]2,\zeta_{18})$ carries them to standard qutrit coordinates $(X,Y,Z)$. On
$\mathfrak c$ the complementary Pfaffian is
\[
  2^{17}C_3N_9S_9^3,
\]
where $C_3N_9$ is, up to $-27$, the product of the twelve stabilizer-state mirrors and
$S_9$ is the product of the nine SIC mirrors. Consequently the Pfaffian weights the two
mirror orbits by $1$ and $3$; the $G_{26}$ Jacobian weights them by $2$ and $1$.
\cert{w33\_pass11261\_g26\_cartan\_coordinate\_map.py}
\end{result}

The two relative invariants separate the two mirror orbits arithmetically. Let $A$ be
the product of the nine SIC mirror forms and $B$ the product of the twelve
stabilizer/MUB mirror forms. Their divisor valuations are
\[
 \operatorname{div}P=(3,1),\qquad \operatorname{div}J=(1,2),
\]
so the valuation matrix has determinant $5$ and
\[
 2\operatorname{div}P-\operatorname{div}J=(5,0),\qquad
 3\operatorname{div}J-\operatorname{div}P=(0,5).
\]
Equivalently $P^2/J$ is a scalar multiple of $A^5$ and $J^3/P$ a scalar
multiple of $B^5$; with the frozen normalisations these become exact cleared
polynomial identities of degrees $78$ and $99$. Thus the signed-$E_8$ Pfaffian
and the $G_{26}$ Jacobian form two mixed probes whose integer combinations
isolate the SIC and stabilizer wall arrangements separately. The index five is
only the determinant of this divisor lattice; it is not a physical $\Z_5$,
particle count or coupling. \tierC\ (\cert{w33\_20261001\_g26\_e8\_mirror\_separation.py})

The same exact Cartan plane admits a reconstructed positive grade-pair form from
the numerical $E_8$ bracket coordinates. In the committed basis,
\[
G_{ij}=\mathrm{Tr}\!\left(\mathrm{ad}(q_i\in\mathfrak g_1)\,
                              \mathrm{ad}(q_j\in\mathfrak g_2)\right)
 =\begin{pmatrix}590&20&0\\20&980&0\\0&0&300\end{pmatrix},
\]
whose principal minors are $590$, $577800$ and $173340000$. Thus, after the natural
$\mathfrak g_1\leftrightarrow\mathfrak g_2$ grade-exchange identification, the semisimple slice has an
positive kinetic candidate for this declared identification. This choice is
additional structure and is not identified with the unitary $G_{26}$ metric.
Using it without fitted coefficients gives a useful failure mode. If
$F=C_3N_9S_9$ is the product of all $21$ mirror forms, the projective functional
$|F|^2/(q^\dagger Gq)^{21}$ has an exact strict local maximum at $[0{:}0{:}1]$:
both real and imaginary tangent Hessians of its logarithm are negative definite
conditional on the displayed reconstructed Gram matrix.
That ray maps to the qutrit $[1,\zeta_9,\zeta_9^{-1}]$ and its cubic postselection
interface has zero Jarlskog invariant. Equal mirror repulsion therefore does
\emph{not} generate CP; an additional oriented datum is required.
\tierC\ (\cert{w33\_20261001\_cartan\_kinetic\_wall\_vacuum.py})

The qutrit ray just displayed is exactly the standard qutrit $T$-magic ray
$|T_3\rangle=(1,\zeta_9,\zeta_9^{-1})/\sqrt3$. For the coefficient-free
equal-wall logarithmic barrier in the standard unitary $G_{26}$ metric it is an
exact stationary point whose real projective Hessian has spectrum
$30,30,54,54$. Its projective Clifford orbit has $72$ rays, stabiliser $C_3$,
and every ray singles out exactly one of the four stabilizer MUBs by the purity
pattern $(1/3,5/9,5/9,5/9)$, with $18$ rays per MUB. Through the dictionary
below this selects one tetrahedral $\phi$ axis; the residual $C_3$ cycles the
other three axes, so one further ternary selector chooses an incident oriented
chord, all of which lie in the TM$_1$ class of Pass~11267. Complex conjugation
stays inside the same $72$-ray orbit and leaves
$(u_6,u_{12},u_{18})=(0,0,-1/729)$ unchanged, so no scalar $G_{26}$-invariant
Cartan potential can choose the CP/time orientation.
\tierC\ (\cert{w33\_20261001\_balanced\_g26\_tmagic\_vacuum.py})

There is also a regular polynomial selector with an exact global statement.
For $\lambda,\alpha,\beta,r_0>0$,
\[
 V(v)=\lambda(v^\dagger v-r_0^2)^2+\alpha|u_6(v)|^2+\beta|u_{12}(v)|^2
\]
has precisely the $72$ projective $T$-orbit rays as global minima at radius
$r_0$, with a common phase still free. In the chart $z=1$, the cube ratios
$A=x^3$, $B=y^3$ have eight simple solutions; the exact eliminant factors as
$(A^2+A+1)(A^3-3A^2-24A-1)(A^3+24A^2+3A-1)$.
All coordinates are nonzero and there are no points at infinity, so the nine
cube-root lifts give $72$ reduced rays. The invariant quotient identifies
these with the $T$ orbit. Their location is independent of positive angular
coefficients. In standard real kinetic coordinates the Hessian spectrum is
$0$, $8\lambda r_0^2$, $32\alpha r_0^{10}$ (twice) and
$200\beta r_0^{22}$ (twice). This is a constructed scalar effective potential;
its higher-degree interactions, radius, couplings and phase require physical
input. Neither these curvatures nor its chosen zero energy predict masses or
the cosmological constant.
\tierC\ (\cert{w33\_20261001\_g26\_global\_polynomial\_vacuum.py})

An alternative normalization can now be certified without reconstructing the
large numerical gauge map. In the canonical flat $27\times3$ order, let $Q$
have columns supported at $(0,52,80)$, $(5,37,75)$ and $(26,40,54)$, with entries
one. The integer $E_6$ basis gives $Q^\dagger Q=3I$, vanishing cross moment
maps and pairwise zero brackets. Thus $\Phi=Qq/\sqrt3$ has canonical norm
$\|q\|^2$ and is exactly $D$-flat. Its mass pencil splits into three cycle
blocks of size three, nine bipartite blocks of size six and one of size eighteen.
Relative signed permutations have order three (up to a minus sign in the
six-blocks); the eighteen-block carries the regular $\Z_3^2$ character.
Consequently, for every complex $q$, its squared singular spectrum is the
twelve normalized MUB overlaps twice each, the nine normalized SIC overlaps
with factor $2/3$ six times each, and three zeros. In particular the moments
are $20\|q\|^2$ and $8\|q\|^4$. This certifies the earlier numerical mirror
law on a new exact plane, not the earlier floating gauge transformation.
\tierC\ (\cert{w33\_20261001\_exact\_cartan\_mirror\_completion.py})

There is an explicit mass mechanism on this normal slice. The signed $E_6$
cubic supplies a global degree-six contraction $C_6$ of three quadratic cross
products; the ternary cubic $d(\Phi a)^3$ supplies an Aronhold scalar $S$.
Exact symbolic restriction gives $C_6|_Q=12u_6$ and
$S|_Q=(152u_6^2-32u_{12})/5$. Thus
$I_6=9C_6/4$ and $I_{12}=(152I_6^2-3645S)/32$ restrict to the standard
invariants on $Q/\sqrt3$. All lower/dual $E_6$ and volume-form generator
checks vanish exactly, and their global analytic differentials are executable. At unit $T$, the covectors
$\overline T,\nabla u_6,\nabla u_{12}$ have Gram $\operatorname{diag}(1,16,100)$.
Their rank-one mass operators therefore lift all three null modes without
changing the $78$ heavy singular masses. With specified EFT coefficients,
the light masses are $|c_5|r^2/\Lambda$, $16|c_6|r^{10}/\Lambda^9$ and
$100|c_{12}|r^{22}/\Lambda^{21}$. The global covariant circuits are explicit;
measured scales, particle assignments and complete EFT-loop stability remain
open. These are nonsupersymmetric EFT interactions.
\tierC\ (\cert{w33\_20261001\_global\_e6\_cartan\_covariants.py})
Adding positive moment squares to the global radial/invariant selector names
an $81$-field classical potential. Its zero set reduces to the earlier
Cartan intersection by moment-zero orbit theory. The $72$ Cartan rays are
little-Weyl-related gauge representatives, not $72$ distinct physical vacua
or prepared quantum resources; the common phase remains flat.

A separate declared gauge/scalar/fermion model does not automatically select
this vacuum. With $g_{SU(3)}=g_{E_6}/\sqrt6$, exact rational inverse-Gram
identities match the gauge squared masses to one third of the cubic squared
masses. The moment potential supplies matching massive real scalars, and
conjugate fermion multiplets cancel perturbative gauge anomalies. The complete
Landau-gauge one-loop angular potential is
$(4g^4/9-8|y|^4)F/(64\pi^2)$, where the earlier interval witnesses show $T$
is a saddle of $F$. A zero prefactor leaves angular flatness. Higher EFT
operators, additional sectors and stable quantum alignment remain open.
\tierC\ (\cert{w33\_20261001\_complete\_gauge\_scalar\_fermion\_loop.py})


The invariant compiler now also supplies the global degree-eighteen coordinate.
If $\mathcal T=dS_F[\operatorname{Hess}(F)]$ is the classical Hessian
polarization of the ternary cubic's Aronhold scalar, exact restriction gives
\[
 I_{18}=\frac{19683\mathcal T-(48384/5)I_6^3+(13824/5)I_6I_{12}}{1492992},
 \qquad I_{18}(Qq/\sqrt3)=u_{18}(q).
\]
Its analytic differential is executable on all $81$ fields. The real term
$\gamma|I_{18}+r_0^{18}/729|^2$ lifts the common phase, but all selected
representatives share one gauge orbit and do not choose physical CP.
\tierC\ (\cert{w33\_20261001\_degree18\_phase\_completion.py})

Actual charges impose a stronger physical boundary. Eight exact stabilizer
generators annihilate the entire $Q$ plane, so the three completed light
modes are singlets of the connected unbroken gauge group. Its dimension eight
cannot contain an unbroken Standard Model algebra of dimension twelve.
The declared conjugate-paired matter has zero chiral index under every
subgroup restriction. Standard $27$ branching does not identify these normal
modes as charged families.
\tierC\ (\cert{w33\_20261001\_chiral\_mass\_assignment\_audit.py})
Including the scalar selector and all three lifts in the full declared
one-loop determinant yields $83$ positive real scalar modes and $81$ massive
fermion singular pairs at $T$. The tested finite-step quantum curvature is
negative and infrared-sensitive because of vanishing Goldstone masses; no
smooth quantum stability theorem follows.
\tierC\ (\cert{w33\_20261001\_completed\_eft\_one\_loop.py})
The same inventory gives standard one-loop $b_{E_6}=17$ and
$b_{SU(3)}=-59/2$, so equal standard couplings drift and family $SU(3)$ has a
Landau pole. Neither transmutation nor a finite vacuum counterterm is fixed.
\tierC\ (\cert{w33\_20261001\_uv\_running\_scale\_audit.py})

There is a distinct conditional nonperturbative model. The $E_6$-only
stabilizer has dimension fourteen on the generic SIC wall $(1,1,z)$ over
$\Q(z)$, but remains eight at the MUB-only control $(0,1,2)$.
A degree-eighteen invariant condensate denominator vanishing on the SIC
orbit is therefore proportional to $u_{18}$, using its order-two mirror
multiplicities. This provides a normalization interface to the near-SUSY
three-$27$ condensate studied at \url{https://arxiv.org/abs/2505.07931}; it
does not import that theory's dynamics into the preceding paired EFT.
For the declared canonical Cartan action
\[
 W=\Lambda^9(-729I_{18})^{-1/3},\qquad
 V=\sum_i|\partial_iW|^2-18m\operatorname{Re}W,
\]
the chosen branch has a stationary radius $r^8=14\Lambda^9/(3m)$.
All six real slice gradients vanish exactly, and the canonical Hessian
has eigenvalues
\[
 m^2\left(\frac{162}{49},\frac{162}{49},\frac{486}{49},
 \frac{486}{49},\frac{648}{7},\frac{486}{7}\right).
\]
The invariant-gradient angular basis has nonzero two-by-two mixing blocks;
these are scalar modes, not CKM or PMNS angles. This result is conditional
on the denominator and canonical slice K\"ahler metric. Full eleven-modulus
stability, quantum K\"ahler corrections, the AMSB/strong scales, Standard Model
charges and nonzero vacuum energy remain separate questions.
\tierC\ (\cert{w33\_20261001\_condensate\_cartan\_potential.py})

\paragraph{Full-modulus and Standard Model interfaces (Passes11270--11274).}
The signed81-field action has an $E_6$ orbit of complex dimension70, so its
physical tangent quotient has eleven complex directions. A numerical full
Hessian and fermion calculation reproduces the published canonical spectrum
of Goh et al.: eight scalar Goldstones, ten squared masses $162m^2/49$,
two $486m^2/49$, and $486m^2/7,648m^2/7$. Two small radial
K\"ahler deformations are tested with the full compensator potential, rather
than an inverse-metric-only substitution; general corrections remain open.
In a separate Higgs configuration, the actual canonical vectors $e_0,e_1$
and an adjoint with Cartan coefficients $(1/3,2/3,1,0,1/2,0)$ leave exactly
$\mathfrak{su}(3)\oplus\mathfrak{su}(2)\oplus\mathfrak u(1)$ with the
standard hypercharges. The signed cubic $d_{AB0}$ is an invariant rank-ten
exotic mass block. Symmetric family-sextet coupling gives rank30 across
three27s, preserving three chiral16s and three singlets. A positive compact
orbit-distance potential selects this imposed Higgs orbit; its extra fields
and nonpolynomial EFT are not a dynamically derived renormalizable vacuum.
\tierC\ (\cert{w33\_pass11270\_full\_condensate\_moduli.py},
\cert{w33\_pass11271\_canonical\_sm\_higgs\_bridge.py})

In the canonical eleven-modulus EFT a computed hard one-loop radial tadpole
gives $\Delta_{\rm hard}=V'_{1,\rm hard}/(2r)$; at the shifted stationary
point $G_{\rm tree}+\Delta_{\rm hard}=0$ for the eight global Goldstones.
Heavy gauge thresholds and momentum-dependent nonzero poles are excluded.
For the supplied torus benchmark, uniform Gaussian-polynomial domination
now proves the existing Wilson curvature response converges at fixed heat
time. At $t=1/5$ the directed outside-$K=64$ tail is below
$1.083\times10^{-25}$; no spacetime emergence or finite-spacing error
enclosure is implied. Finally the canonical spectrum obeys
$\operatorname{STr}M^2=0$ but
$\operatorname{STr}M^4=(852930/2401)m^4>0$, retaining vacuum-energy
running. An external global-sequestering action passes the constant-shift
test, but its constraints and historic residual are not supplied by W33.
\tierC\ (\cert{w33\_pass11272\_resummed\_condensate\_radial\_control.py},
\cert{w33\_pass11273\_wilson\_uniform\_tail.py},
\cert{w33\_pass11274\_scale\_and\_sequestering.py})

\paragraph{Polynomial local alignment and new physical tests (Passes11275--11279).}
A separate polynomial Higgs EFT on two complex27 vectors and a real78
has a nonnegative sum-of-squares potential whose exact constraint differential
has rank120 in186 real fields. Its kernel is exactly the66-dimensional compact
$E_6$ orbit of the canonical SM point, so the orbit is isolated locally with
positive normal Hessian. This uses an imposed degree8 cubic-covariant selector;
global uniqueness and selected scales remain open. Two added complex matrix
scalars give a degree-four sum-of-squares lift with the same local gauge kernel,
but their large inventory changes the combined $E_6$ beta coefficient to$-80$.
This renormalizable scalar construction is not an established viable UV model.
An independent family sextet produces a rank-ten spectator mass through a
dimension-six composite operator, improving the family beta coefficient to
$-79/3$ without restoring asymptotic freedom. An actual momentum bubble in a
\emph{separately declared radial-only scalar EFT} shifts its squared pole by
about $-0.144\%$ from curvature; other modulus, fermion and gauge loops are
excluded. A mixed-characteristic symplectic lift is explicit, but does not
select an invariant positive real metric. The known clique topology also makes
the top boundary of its circle-history product injective: this candidate has
no intrinsic topological four-flux for sequestering. Neither geometry test
excludes all possible W33-derived spacetime constructions.
\tierC\ (\cert{w33\_pass11275\_polynomial\_sm\_higgs.py},
\cert{w33\_pass11276\_sextet\_composite\_spectator.py},
\cert{w33\_pass11277\_radial\_momentum\_pole.py},
\cert{w33\_pass11278\_symplectic\_ring\_geometry.py},
\cert{w33\_pass11279\_history\_flux\_obstruction.py})

\paragraph{Smaller field and continuum architectures (Passes11280--11284).}
A degree-four factor lift of the rank-five cubic selector uses three complex
$27\times5$ fields with an added auxiliary $U(5)$ gauge factor. Its996 real
fields have91 gauge and905 positive normal directions locally; the combined
$E_6$ coefficient becomes13 and auxiliary $SU(5)$ coefficient$29/6$, while
the auxiliary Abelian and existing family $SU(3)$ UV problems remain. A degree-twelve family-sextet
potential selects an imposed orbit with singular values$(1,2,3)$; a single
sextet tree renormalizable potential cannot isolate three distinct nonzero
values. The mixed bilinear $v^A H^\dagger_A$ avoids the earlier neutral-Higgs
cubic zero, but one aligned texture gives no CKM mixing.
The radial self-energy now includes all22 scalar and11 Weyl internal modes
in the declared fixed-slice EFT: its open Goldstone channels give
$\Gamma(h\to8GG)=M_h^3/(8\pi r_0^2)$ and a complex pole with leading
$\Gamma/M_h\simeq0.003606$ for the stated parameters. The full external
self-energy matrix and physical UV matching remain open.
An added real continuum metric chart needs a relative symplectic-plane scale
and an overall scale to span all ten local Lorentz-metric variations. An added
identity double of the genus81 Levi-graph handlebody, times a history circle,
has one top-flux class and supports a conditional linear-sequestering action.
These are explicit architectures with additional real-base, gluing, coupling
and flux inputs; they do not derive gravity or the observed vacuum energy.
\tierC\ (\cert{w33\_pass11280\_factor\_higgs\_mediator.py},
\cert{w33\_pass11281\_family\_sextet\_selection.py},
\cert{w33\_pass11282\_all\_modulus\_radial\_self\_energy.py},
\cert{w33\_pass11283\_symplectic\_metric\_field.py},
\cert{w33\_pass11284\_closed\_history\_flux.py})

A further semisimple auxiliary SO(10) completion needs explicit frame
projection constraints to remove quartic-flat unused columns. Its local normal
count is1740; replacing the elementary family Higgs by a composite leaves
$b_{E_6}=4$, $b_{SO(10)}=1$, with UV mediators and family running still open.
A separately imposed CP-even mixed-moment potential selects Fourier mixing
with $J^2=1/108$, rather than the observed CKM hierarchy. The canonical
K\"ahler quotient now supplies the full22-field nonanalytic self-energy matrix,
checked by an independent nonradial evaluation on the actual D-flat sheet;
local matching, vector thresholds and physical poles remain open. The imposed
Einstein--Hilbert metric pullback has two graviton polarizations under its
fiber and ADM constraints. Nontrivial Levi-cycle gluing preserves one top-flux
class but does not select its magnitude. The cycle determinant and critical
group are prior results (Passes5031/5441); their separately declared Abelian
Chern--Simons model has a gapped-boundary obstruction repaired by doubling.
These constructions retain additional potential, geometry and coupling inputs.
\tierC\ (Passes11285--11292; scoped certificates and independent controls)

Explicit tree matching of the composite Yukawa costs six or twelve $E_6$
beta units, exceeding the old budget of four. Replacing the third large frame
by a real auxiliary symmetric54 restores a quartic alignment potential with
1254 positive normal directions; the two mediator choices retain positive
$(b_{E_6},b_{SO(10)})=(8,8)$ or $(2,8)$. Family running and complete-coupling
UV stability remain open. A separately declared hierarchical CP potential
uses adjugate-trace invariants to remove the four relative sextet flats at four
tested vacua; its hierarchy and root coefficients remain inputs. Local two-point
Ward matching yields fourteen leading complex massive poles plus eight
Goldstones in a declared spacelike scheme, with invariant tadpole and UV
matching still open. Minimal covariant quotient matter preserves the metric
fiber and continuum ADM constraints. An added membrane-rate model selects a
nearest-lattice flux sector but retains bare vacuum energy. The old uniform
linear-volume branch fixes the intensive $Q/\mathrm{Vol}$, so its extensive
flux is not that membrane ladder. These are scoped EFT constructions and
obstructions, not observed-parameter predictions.
\tierC\ (Passes11293--11297; matching, phase, constraint and flux controls)

The next audit contracts the actual single-channel mediator Yukawas, but
finds that one vectorlike pair has a rank-one channel map and cannot implement
two independent flavour sectors. A second pair gives $b_{E_6}=-10$; two scalar
mediators instead retain $b_{E_6}=2$, with their complete running open.
Mandatory symmetric54 quartics expose missing counterterms; full portal
closure remains open. For the canonical two-Higgs mass map, five signed paths
prove that its entire spectrum depends only on $S_u+S_d$, so this fermion
determinant cannot lift relative Takagi phases or derive the imposed hierarchy.
An invariant Hermite counterfunctional matches the vector-sector tadpole and
Hessian in a declared higher-degree EFT; scalar/Weyl and UV matching remain
open. Native Levi site constraints generate missing edge currents, rather
than inheriting continuum ADM closure. A separately added Maxwell membrane
form admits a closed two-cap junction solution with volume and curvature-flux
constraints and common bare-shift cancellation. Its geometry, flux sectors
and scales remain inputs; no observed-parameter prediction is claimed.
\tierC\ (Passes11298--11302; scoped tensor, matrix, bracket and junction controls)

A single scalar mediator with two sources escapes the fermion-pair channel-rank
obstruction. An exact Riccati inequality excludes one-loop complete weak-coupling
asymptotic freedom of the unchanged54 inventory including scalar portals; eight
added SO10 vector Weyls, two coupled to54, supply a bounded AF ray in the isolated
scalar/fermion block. Full portal completion and exotic-field decoupling remain
open. Two vectorlike E6-singlet family probes detect all four old Takagi-phase
flats, without deriving a stationary physical vacuum. Full vector cut matrices
extend conditional massive poles, but a rank11 scalar Goldstone infrared
logarithm prevents unqualified regulator-free total local matching. Added native
Levi clocks admit exactly first-class relational constraints, retaining80
oscillators rather than deriving gravity. In a separately supplied two-cap
membrane model, integer flux leaves one constrained radial negative direction
with curvature about$-204.71$; full fluctuation stability, geometry and observed
cosmological-constant selection remain open.
\tierC\ (Passes11303--11307; scoped EFT, tensor and constrained-action controls)

A subsequent complete seven-quartic $54+45$ restriction exposes a further
obstruction: the original eight-Weyl texture has a positive $45$ Riccati
barrier, of minimum $298418/295695$, including its scalar portals. An allowed
joint $54/45$ Yukawa texture instead has a positive $54$ barrier. The earlier
isolated AF ray remains valid; neither result closes the full inventory.
Its isolated $54$ model admits a positive leading radiative radial curvature
on an input RG crossing, but selects an $\mathrm{SO}(9)$ orbit and leaves the
absolute scale and vacuum energy free. Untargeted flavor loops yield a
CP-conserving local five-phase minimum, with spectra and angles supplied.
The massless Goldstone bubble has a regulator-free limit at massive external
momentum; invariant tadpole and UV matching remain open. A constructed native
edge Fierz--Pauli action on a supplied four-dimensional base carries two uniform
massless polarizations and $79$ massive spin-two modes, without a nonlinear
completion. The declared identity-gluing history has $\chi=0$, $b_1=82$ and
therefore no smooth closed Riemannian Einstein metric: Einstein plus $\chi=0$
would force a flat finite torus cover, whose $b_1$ is at most four. This does
not exclude walls, matter, Lorentzian histories or other gluings. A constructed
local smooth neck requires radial null-energy violation and is not a global
membrane saddle.
\tierC\ (Passes11313--11319; exact obstructions and conditional EFT constructions)

Extending the real two-active-flavour Yukawa family gives eight exact support
patterns, including an unequal noncommuting branch, but all retain a positive
scalar AF barrier. A supplied full24-field family-invariant EFT has numerical
local minima with spontaneous relative CP and $2+1$ singular-value splitting;
commuting Gram matrices still preclude a CKM prediction. Finite invariant
hard-sector spectral counterfunctions give conditional invariant hard tadpole
matching; nonlinear covariant quotient maps, the
rank11 soft logarithm and complete resummation remain open. Actual native-graph
shift elimination in an added nonlinear pair-metric action gives a rank78 lapse
Hessian; no symmetry-invariant spanning tree exists. A distinct, chosen
$S^2\times T^2$ history admits an exact positive-tension Maxwell-supported wall
satisfying the bulk Einstein equations, both wall junctions and magnetic Dirac
quantisation. An explicit native cochain selects one primitive handle; this is
not the full81-handle history. Its coupled fluctuation spectrum, state selection
and observed vacuum energy remain open.
\tierC\ (Passes11320--11324; scoped EFT witnesses and exact constraint/geometry tests)

Complex two-active Yukawa fixed rays phase-lock to the prior real families;
paired two/four/six/eight-flavour rays retain positive scalar AF barriers.
A supplied bounded CP-even angular interaction yields a local noncommuting
24-field flavour minimum with sixteen positive normal Hessian eigenvalues
and Jarlskog $0.09622$. Its nearly maximal mixing and strong engineered EFT
coupling are not an observed CKM prediction. An explicit numerical local
22-coordinate D-flat quotient now supplies the metric, connection and covariant
scalar/Weyl/vector mass maps missing in Pass11322. Reference spectra,
changed-frame covariance and quadratic remainders against the old linear jets
are checked; full hard 1PI Ward resummation remains open.
The published determinant multivielbein construction has a restricted equal-boost
constraint proof. Applied to eighty vielbeins it gives 397 modes on that branch,
but a single vertex couples 3160 pairs rather than the native 160 edges.
On the actual Maxwell-wall saddle, both torus-shape metric polarizations have
positive Euclidean Rayleigh forms and two exact constant zero modes.
The remaining coupled fluctuation block, a generic native nonlinear gravity
completion and observed physical scales remain open.
\tierC\ (Passes11325--11329; scoped local witnesses and algebraic tests)

An exact complex-isotropic four-flavour Yukawa ray evades the old scalar
lower-bound tests, but an actual scalar background with fourth trace $1/5$
gives a new positive Riccati minimum $372423967259/182902583664$ and excludes
simultaneous seven-quartic AF fixed points on this branch. Gaussian positive-mass
adjoint mediation instead forces commuting flavour Grams. Compensated family
vector fields on the local quotient fix the eight hard Ward columns; the
normal hard1PI block remains uncomputed. On a collinear native pair-vielbein
slice, relative boosts are fixed independently of lapses and the reduced action
is lapse-linear despite cycles; generic transverse constraints remain open.
The coupled Maxwell/metric torus-vector block of the original wall has four
exact zero modes, hence six with the previous shapes. A supplied two-axion
winding extension admits a backreacted quantised wall with $q^2-q=h b_{\rm wall}$
and lifts both shapes: its smallest shape eigenvalue lies in $[8/625,1/50]$.
The remaining coupled sectors, nonlinear integrability, matter origin and
physical scale or vacuum-energy selection remain open.
\tierC\ (Passes11335--11341; scoped exact and numerical constructions)

Both stored four-active complex Yukawa witnesses that escaped the older scalar
portal tests fail the stronger physical-background bound. A bounded, positive-mass
two-adjoint mediator sector selects noncommuting two-family Grams at fixed spectra,
but gives no three-family CP or RG completion. The eight quotient Ward columns
leave 105 independent entries of the normal hard Hessian unfixed; the native
graph's three transverse boost directions have local shift rank 237, without
establishing nonlinear Dirac closure. On the backreacted axion-winding wall,
$\chi=1/b-4/5$ and $V_\phi=-F/(qcb^2)$ solve the coupled vector zero equations
exactly, even though winding lifts the two torus-shape zeros. The same exact
family has a positive, quantised $\rho=0$ bare-cosmological-parameter point at
$h=16/45$, $q=4/3$, with supplied matter strength and wall tension. Neither
construction fixes an observed mass, absolute scale or vacuum energy.
\tierC\ (Passes11342--11349; exact restricted identities and finite controls)





The computational interface needs an additional classical datum: a Clifford
frame for the selected ray. The $72$ rays split into eight nine-ray Pauli
orbits, two per MUB axis. Uniformly forgetting the displacement within even
one such orbit gives $I/3$, by the symbolic Pauli-twirl identity. Retaining a
representative $C$ permits $C^\dagger$ followed by \textsc{sum} with a fresh
$|0\rangle$ to prepare $(I\otimes T)|\Phi_3\rangle$. All $72\times9=648$
ray/Bell branches replay the existing injection identity
$C_{ab}K_{ab}=T/3$. If the frame is uniformly forgotten, the same protocol
instead gives the entanglement-breaking channel $\rho\mapsto\operatorname{diag}\rho$.
Thus orbit selection and frame retention are separate resources; a
state-preparation Hamiltonian and its noise model remain open.
\tierC\ (\cert{w33\_20261001\_tmagic\_orbit\_factory\_contract.py})

The twelve stabilizer mirrors form the four qutrit mutually unbiased bases. Complex-linear
$G_{26}$ reflections induce $A_4$ on those four bases; complex conjugation supplies an odd
transposition and completes the image to $S_4$. This identifies the line-stabiliser
$S_4$ flavour tetrahedron objectwise: four flavon axes are four MUBs and the twelve
oriented tetrahedral chords are the ordered MUB pairs. It does not fix Yukawa amplitudes;
in particular a vacuum-tilt parameter does not determine $\theta_{13}$ or $\delta$ by
itself. \cert{w33\_pass11262\_g26\_mub\_s4\_yukawa\_bridge.py}

There is also a metric-independent spectral statement. The nonzero spectrum of $N^3$ is
the same in all three grades, so the $234$ nonzero modes form $78$ exact $\Z_3$ triplets.
Every positive-power root-of-unity graded trace cancels; the residual graded index is
$8+3\omega+3\omega^2=5$ and lies wholly in the kernel. This is a character identity,
not a boson--fermion supertrace or a vacuum-energy calculation.
\cert{w33\_pass11263\_full\_graded\_spectrum.py}

\subsection{Clock gradings of \texorpdfstring{$E_8$}{E8}}

A clock gives more than a $\Z_3$-grading: once an origin and an orientation are chosen it
gives an \emph{integer} grading. On the committed Chevalley basis:

\begin{result}{Theorem 5.5 (Hesse clocks grade $E_8$) \hfill \tierC}
\begin{enumerate}[label=(\roman*),leftmargin=1.5em]
\item A single Hesse line (a ``tick'') induces the \emph{contact grading}
      \[ E_8=\mathfrak g_{-2}\oplus\mathfrak g_{-1}\oplus\mathfrak g_0\oplus\mathfrak g_1\oplus\mathfrak g_2,\qquad 1+56+134+56+1, \]
      with $\mathfrak g_0=E_7\oplus\C$.
\item An oriented striation (three parallel ticks) induces the \emph{$|3|$-grading}
      \[ 2+27+54+82+54+27+2, \]
      with Levi $\mathfrak g_0=E_6\oplus A_1\oplus\C$, $\mathfrak g_{\pm1}=\mathbf{27}\otimes\mathbf2$,
      $\mathfrak g_{\pm2}=\mathbf{27}$, $\mathfrak g_{\pm3}=\mathbf2$.
\item The repository's frozen $\Z_3$-grading $86+81+81$ is exactly this integer grading
      reduced modulo three; all $8\,347$ nonzero committed structure constants obey the
      integer degree law.
\end{enumerate}
\cert{w33\_20260925\_hesse\_clock\_e8\_gradings.py},
\cert{w33\_20260925\_e8\_parabolic\_cubic\_clock\_lift.py}
\end{result}

\begin{figure}[htbp]
\centering
\begin{tikzpicture}[yscale=0.028,xscale=1]
  % contact grading
  \begin{scope}
    \foreach \d/\h/\c in {-2/1/ink,-1/56/teal,0/134/gold,1/56/teal,2/1/ink}{
      \fill[\c!70] (\d*0.9-0.32,0) rectangle (\d*0.9+0.32,\h);
      \node[font=\scriptsize,anchor=south] at (\d*0.9,\h+2) {$\h$};
      \node[font=\scriptsize,anchor=north] at (\d*0.9,-3) {$\d$};}
    \node[font=\sffamily\scriptsize,text=slate] at (0,-22) {contact grading (one tick)};
    \node[font=\scriptsize,text=white] at (0,40) {$E_7$};
    \node[font=\scriptsize,text=white] at (0,28) {$+\C$};
    \node[font=\tiny,text=white,rotate=90] at (-0.9,26) {Freudenthal};
    \node[font=\tiny,text=white,rotate=90] at (0.9,26) {Freudenthal};
  \end{scope}
  % |3| grading
  \begin{scope}[xshift=7.2cm]
    \foreach \d/\h/\c in {-3/2/ink,-2/27/rose,-1/54/teal,0/82/gold,1/54/teal,2/27/rose,3/2/ink}{
      \fill[\c!70] (\d*0.9-0.32,0) rectangle (\d*0.9+0.32,\h);
      \node[font=\scriptsize,anchor=south] at (\d*0.9,\h+2) {$\h$};
      \node[font=\scriptsize,anchor=north] at (\d*0.9,-3) {$\d$};}
    \node[font=\sffamily\scriptsize,text=slate] at (0,-22) {$|3|$-grading (oriented striation)};
    \node[font=\tiny,text=white,align=center] at (0,33) {$E_6$\\$A_1$\\$\C$};
  \end{scope}
\end{tikzpicture}
\caption{Two ways a clock slices $E_8$. Heights are dimensions. On the left, the
Freudenthal ``$56$'' of Section~6 appears at degree $\pm1$; on the right, the $27$ of
$E_6$ appears doubled at degree $\pm1$ and single at degree $\pm2$.}
\label{fig:gradings}
\end{figure}

This parabolic spectrum was found independently by Kraft, Regeta and Zimmermann in a
study of small varieties with group actions \cite{KRZ2020}. What is repository-specific is
the identification with the Hesse clocks and the committed structure constants.

\subsection{The cubic ladder: three steps to the top}

In the $|3|$-grading, positive degrees form a nilpotent algebra
$\mathfrak g_+=\mathfrak g_1\oplus\mathfrak g_2\oplus\mathfrak g_3$ of dimension $54+27+2=83$. It is generated by $\mathfrak g_1$, and
the way it fills up is governed by the $E_6$ cubic.

\begin{result}{Theorem 5.6 (the bracket ladder) \hfill \tierC}
On the committed Chevalley bracket:
\begin{enumerate}[label=(\roman*),leftmargin=1.5em]
\item $[\mathfrak g_1,\mathfrak g_1]=\mathfrak g_2$: $270$ nonzero pairs hit all $27$ grade-two roots, each ten times;
\item $[\mathfrak g_1,\mathfrak g_2]=\mathfrak g_3$: $54$ nonzero pairs hit both top roots, each $27$ times;
\item the nonzero nested triple brackets exhibit exactly the $45$ triads of the cubic,
      each $12$ times, and for every signed triad $d_{ijk}$,
      $[[e_{(i,0)},e_{(j,1)}],e_{(k,a)}]=-d_{ijk}\,T_a$ for $a\in\{0,1\}$.
\end{enumerate}
So the positive part grows $54\to81\to83$, and the committed cubic is exactly the rule by
which three first-order moves reach the two-dimensional top.
\cert{w33\_20260925\_cubic\_clock\_bracket\_ladder.py}
\end{result}

\begin{figure}[htbp]
\centering
\begin{tikzpicture}[>={Stealth[length=2.3mm]},
   blk/.style={draw=#1,fill=#1!12,rounded corners=3pt,align=center,font=\sffamily\scriptsize,minimum height=9mm}]
  \node[blk=teal,minimum width=36mm] (g1) at (0,0) {$\mathfrak g_1=\mathbf{27}\otimes\mathbf2$\\$54$ first-order moves};
  \node[blk=rose,minimum width=26mm] (g2) at (4.6,0) {$\mathfrak g_2=\mathbf{27}$\\$27$ histories};
  \node[blk=ink,minimum width=18mm] (g3) at (8.3,0) {$\mathfrak g_3=\mathbf2$\\top doublet};
  \draw[->,line width=1pt] (g1) -- node[above,font=\scriptsize]{$[\mathfrak g_1,\mathfrak g_1]$} node[below,font=\tiny,text=slate]{$270$ pairs, $\times10$} (g2);
  \draw[->,line width=1pt] (g2) -- node[above,font=\scriptsize]{$[\mathfrak g_1,\cdot\,]$} node[below,font=\tiny,text=slate]{$54$ pairs, $\times27$} (g3);
  \node[font=\sffamily\scriptsize,text=slate] at (0,-1.1) {cumulative $54$};
  \node[font=\sffamily\scriptsize,text=slate] at (4.6,-1.1) {cumulative $81$};
  \node[font=\sffamily\scriptsize,text=slate] at (8.3,-1.1) {cumulative $83$};
  \draw[->,gold!80!black,dashed,line width=0.8pt] (g1.north) to[bend left=22] node[above,font=\scriptsize,text=gold!70!black]{three moves reach the top with coefficient $-d_{ijk}$: the $E_6$ cubic} (g3.north);
\end{tikzpicture}
\caption{The cubic ladder. First-order moves combine in pairs to histories, and histories
combine with one more move to reach the two-dimensional top. The only way three moves
reach the top is through a triad of the cubic.}
\label{fig:ladder}
\end{figure}

\begin{physical}[What it means physically: cubic time]
Read degree as ``order in time''. A first-order move is a state update; two moves compose
into a second-order \emph{history}; and three reach the top doublet, whose coefficient is
the cubic volume $N$ spanned by the three moves. In this reading the $E_6$ cubic --- in
GUT language the Yukawa coupling --- is the rule by which three independent updates
produce a net displacement of the clock. The algebra is exact; the reading is a bridge,
and it lacks what would make it physics: a Hamiltonian, a state and an energy scale.
\tierB
\end{physical}

\begin{keyidea}
One qutrit's nine histories build $E_8$ as $80+84+84$. The four Hesse clocks are four
orthogonal $A_2$'s whose glue is the tetracode; a clock tick grades $E_8$ as
$1+56+134+56+1$, an oriented clock as $2+27+54+82+54+27+2$; and $E_8$'s roots fibre
six-to-one over the forty points.
\end{keyidea}
